\documentclass[11pt,a4paper]{article}

% ------------------------------------------------------------------
% House layout, shared with the sibling preprints in docs/research/:
% 11pt a4, 2.5cm margins, Computer Modern, \maketitle front matter,
% abstract + keywords + JEL codes, numbered sections, table of contents.
% ------------------------------------------------------------------
\usepackage{amsmath, amssymb, amsthm}
\usepackage[a4paper, margin=2.5cm]{geometry}
\usepackage{array}

\theoremstyle{definition}
\newtheorem{definition}{Definition}[section]
\newtheorem{remark}{Remark}[section]
\newtheorem{proposition}{Proposition}[section]
\theoremstyle{plain}
\newtheorem{theorem}{Theorem}[section]
\newtheorem{lemma}{Lemma}[section]

\newcommand{\R}{\mathbb{R}}
\newcommand{\E}{\mathbb{E}}
\newcommand{\N}{\mathbb{N}}
\newcommand{\norm}[1]{\left\lVert #1 \right\rVert}
\newcommand{\corr}{\operatorname{corr}}
\newcommand{\sgn}{\operatorname{sign}}
\DeclareMathOperator{\IC}{IC}
\newcommand{\code}[1]{\texttt{\detokenize{#1}}}

% booktabs unavailable in this distribution — plain equivalents
\newcommand{\toprule}{\hline\hline}
\newcommand{\midrule}{\hline}
\newcommand{\bottomrule}{\hline\hline}
% natbib for the author-year bibliography (this paper cites inline)
\usepackage[round]{natbib}

\newcommand{\feat}[1]{\texttt{#1}}
\newcommand{\pending}[1]{\emph{[Pending: #1]}}

\begin{document}

% ============================================================

\title{\textbf{Predictive Structure of the Perpetual-Futures\\[3pt]
Limit Order Book}\\[6pt]
\large Universality, Orthogonality, and the Execution Boundary}

\author{Yigit Onat\thanks{Working paper, preliminary draft --- comments
welcome. The out-of-time replication study
(Section~\ref{sec:replication}) is in progress; all other estimates are
from the internal empirical record consolidated 2026-07-25. Please do
not cite specific numbers without checking for a newer version.}\\
\textit{Independent Researcher}\\
\texttt{yionat@gmail.com}}
\date{July 2026}
\maketitle

% ============================================================
\begin{abstract}
Using a 236-feature representation of the limit order book computed at
100\,ms cadence on BTC, ETH and SOL perpetual futures (2.17 million
ticks per symbol), we document the predictive structure of the book at
horizons of one second to one hundred minutes. Directional information
is large (rank information coefficient up to $0.47$ at 1--5\,s),
universal across symbols and volatility regimes, concentrated in eight
independent feature axes led by order-book imbalance, and spectrally
localized in the 0.005--0.1\,Hz band with an Ornstein--Uhlenbeck
half-life of 5--7\,s. Volatility information (peak IC $0.35$) resides
in a disjoint feature family --- arrival intensity, activity counts,
toxicity --- and the two channels are empirically orthogonal:
directional features carry no volatility information and vice versa,
on all three symbols. The signal decays toward efficiency in real time
(its half-life shortening as spreads tighten) and is sharply bounded
at execution: conditioning on directionally-consistent maker fills
collapses the IC to ${\approx}\,0.03$, while conditioning on the
presence of any fill \emph{raises} it to $0.52$, isolating adverse
selection --- rather than signal decay --- as the binding constraint.
An out-of-time replication on an independent month of data is in
progress. We discuss implications: two-block (sign $\times$ size)
forecasting architectures, entropy-gated conditioning, and
execution-timing overlays for slower strategies.
\end{abstract}

\smallskip
\noindent\textbf{Keywords:} limit order book, market microstructure, cryptocurrency, perpetual
futures, information coefficient, adverse selection, spectral analysis

\smallskip
\noindent\textbf{JEL Classification:} C58, G12, G14

\tableofcontents
\bigskip

% ============================================================
\section{Introduction}
% ==================================================================

That order-book imbalance predicts short-horizon price moves is
well established: \citet{cont2014} document the linear price impact of
order-book events in equities, and a large subsequent literature
confirms the result across venues and with increasingly flexible
models. This paper does not re-claim that finding. Instead, we use a
purpose-built measurement platform --- 236 features computed every
100\,ms from the full depth, trade flow, and funding state of three
cryptocurrency perpetual-futures books --- to establish four facts
that are, to our knowledge, new.

\begin{enumerate}
\item \textbf{Orthogonality (the factorization result).} The
directional and volatility channels of the book are carried by
\emph{disjoint} feature families. Imbalance-family features predict
the sign of returns (IC up to $0.47$) and carry zero volatility
information; intensity-family features predict magnitude (IC up to
$0.35$) and carry zero directional information. This holds feature by
feature, on all three symbols. Sign and size of short-horizon returns
are independently predictable --- a factorization with direct
architectural consequences (Section~\ref{sec:orthogonality}).
\item \textbf{Universality at scale.} A census of 207 live features
across six horizons and three symbols, with bootstrap confidence
intervals, compresses to eight independent directional axes whose IC
ordering is preserved across symbols
(Section~\ref{sec:directional}).
\item \textbf{Spectral localization and real-time efficiency drift.}
The entire directional IC lives in the 0.005--0.1\,Hz band; the
imbalance process has an Ornstein--Uhlenbeck half-life of 5--7\,s; and
between two measurement windows six weeks apart the 5\,s IC decayed
from $0.45$ to $0.29$ while the 1\,s IC held --- the signal's
half-life is shortening as the venue's spreads tighten
(Section~\ref{sec:spectral}).
\item \textbf{The execution boundary.} Conditioning the IC on
directionally-consistent maker fills collapses it from $0.45$ to
${\approx}\,0.03$; conditioning on \emph{any} fill raises it to
$0.52$. The decomposition locates the loss at fill selection, not at
signal decay, and isolates adverse selection in the sense of
\citet{glosten1985} as the binding mechanism
(Section~\ref{sec:boundary}). Regime conditioning on a book-shape
entropy measure recovers part of the boundary (+22\% IC in the
low-entropy quintile).
\end{enumerate}

The perspective throughout is measurement-first: we report what the
book knows, where in frequency it knows it, how the two channels
factorize, and precisely where the information stops being capturable.
Section~\ref{sec:replication} pre-registers an out-of-time replication
of every estimate on an independent month of data.
Section~\ref{sec:implications} discusses implications for forecasting
architecture and execution design.

% ==================================================================
\section{Data and Measurement Platform}\label{sec:data}
% ==================================================================

\paragraph{Venue and instruments.} We study BTC, ETH and SOL perpetual
futures on Hyperliquid, a high-throughput decentralized derivatives
venue. Perpetual futures are the dominant crypto trading instrument
and carry a funding-rate mechanism whose pricing theory is developed
in \citet{ackerer2025}.

\paragraph{Feature computation.} A Rust ingestion engine subscribes to
the full order-book and trade stream and emits a 236-dimensional
feature vector per symbol every 100\,ms, spanning 21 categories: book
imbalance at multiple depths, raw depth, microstructure dynamics
(queue position, imbalance velocity), trade flow and aggressor ratios,
VWAP deviation, entropy measures (permutation entropy, book-shape
entropy), Hawkes-type arrival intensity, realized and range-based
volatility, spread and toxicity proxies, funding/open-interest state,
and cross-symbol imbalance aggregates. End-to-end feature latency is
below 80\,ms at the 99th percentile.

\paragraph{Census window.} The headline census uses three contiguous
days (2026-05-19 to 2026-05-21), 2.17 million ticks per symbol at
100\,ms cadence. Rank (Spearman) ICs against future mid-price returns
are computed at six horizons (1\,s, 5\,s, 30\,s, 1\,m, 5\,m, 15\,m;
selected features also at 100\,m) on ${\sim}21{,}708$ subsampled
points per symbol (one per ${\sim}10$\,s to control overlap), with
significance at $p<0.01$ and multiplicity handled by
Benjamini--Hochberg false-discovery control \citep{benjamini1995}.
Of 236 computed features, 207 were live in the census window.
The fill-conditioning study (Section~\ref{sec:boundary}) uses a longer
window: 31 calendar dates (2026-04-19 to 2026-06-04), of which 23
passed data-quality gates.

\paragraph{Honest accounting.} Two limitations are flagged here rather
than discovered later. First, the census window is short; temporal
out-of-sample stability initially \emph{failed} our own validation
matrix (Section~\ref{sec:robustness}), which is precisely what the
replication study addresses. Second, the platform's data-continuity
record over April--June contains gaps; all estimates are computed on
gap-audited clean days only.

% ==================================================================
\section{The Directional Channel}\label{sec:directional}
% ==================================================================

Table~\ref{tab:directional} reports peak rank ICs for the leading
directional features. Order-book imbalance dominates: at 1--5\,s
horizons the level-1 quantity imbalance reaches IC $+0.45$ to $+0.47$
depending on symbol, with the level-5 and notional variants
essentially tied. The ordering of features is preserved across all
three symbols --- the hierarchy is structural, not asset-specific.

\begin{table}[htbp]
\centering
\small
\caption{Leading directional features: peak rank IC by symbol
(census window 2026-05-19/21; horizons 1--5\,s unless noted).}
\label{tab:directional}
\begin{tabular}{lrrrl}
\toprule
Feature & BTC & ETH & SOL & Axis \\
\midrule
\feat{imbalance\_qty\_l1}        & $+0.453$ & $+0.447$ & $+0.466$ & book imbalance \\
\feat{imbalance\_qty\_l5}        & $+0.456$ & $+0.431$ & $+0.453$ & book imbalance \\
\feat{imbalance\_notional\_l5}   & $+0.456$ & $+0.431$ & $+0.453$ & book imbalance \\
\feat{imbalance\_orders\_l5}     & $+0.435$ & $+0.438$ & $+0.455$ & book imbalance \\
\feat{imbalance\_depth\_weighted}& $+0.434$ & $+0.410$ & $+0.425$ & book imbalance \\
\feat{raw\_bid\_depth\_5}        & $+0.420$ & $+0.405$ & $+0.406$ & depth asymmetry \\
\feat{raw\_ask\_depth\_5}        & $-0.424$ & $-0.404$ & $-0.415$ & depth asymmetry \\
\feat{cross\_obi\_mean}          & $+0.354$ & $+0.334$ & $+0.331$ & cross-symbol \\
\feat{micro\_queue\_position\_bid}& $+0.307$ & $+0.282$ & $+0.291$ & queue dynamics \\
\feat{flow\_vwap\_deviation} (1\,s) & $-0.287$ & $-0.206$ & $-0.188$ & VWAP reversion \\
\feat{micro\_obi\_velocity} (1\,s)  & $+0.192$ & $+0.181$ & $+0.208$ & imbalance momentum \\
\feat{ent\_permutation\_imbalance\_16} (1\,s) & $-0.170$ & $-0.173$ & $-0.177$ & imbalance entropy \\
\feat{flow\_aggressor\_ratio\_5s}   & $+0.112$ & $+0.109$ & $+0.111$ & aggressor flow \\
\bottomrule
\end{tabular}
\end{table}

\paragraph{Eight independent axes.} The ${\sim}29$ fast-directional
features are heavily cross-correlated. They compress into eight
axes with approximate standalone IC:
(1) book imbalance ${\sim}0.45$;
(2) raw depth asymmetry ${\sim}0.42$;
(3) cross-symbol imbalance ${\sim}0.35$;
(4) queue dynamics ${\sim}0.31$;
(5) VWAP deviation, mean-reverting ${\sim}0.25$;
(6) imbalance velocity ${\sim}0.19$;
(7) imbalance entropy ${\sim}0.17$;
(8) aggressor flow ${\sim}0.11$.

\paragraph{Decay.} Table~\ref{tab:decay} traces the peak directional
IC across horizons. The half-life is ${\sim}30$\,s for BTC,
${\sim}20$\,s for ETH and ${\sim}15$\,s for SOL --- the thinner the
book, the faster the decay --- and by five minutes roughly 80\% of the
signal is gone. A distinct \emph{slow} family grows with horizon
instead (e.g.\ a 1-hour divergence measure reaching IC $0.21$ at
100\,m), but its analysis is outside this paper's scope.

\begin{table}[htbp]
\centering
\small
\caption{Peak directional IC by horizon (census window).}
\label{tab:decay}
\begin{tabular}{lrrrrrr}
\toprule
Symbol & 1\,s & 5\,s & 30\,s & 1\,m & 5\,m & 15\,m \\
\midrule
BTC & $0.447$ & $0.456$ & $0.261$ & $0.187$ & $0.090$ & $0.056$ \\
ETH & $0.447$ & $0.438$ & $0.219$ & $0.160$ & $0.074$ & $0.031$ \\
SOL & $0.466$ & $0.412$ & $0.190$ & $0.135$ & $0.056$ & $0.036$ \\
\bottomrule
\end{tabular}
\end{table}

% ==================================================================
\section{The Volatility Channel}\label{sec:volatility}
% ==================================================================

Table~\ref{tab:volatility} reports the leading predictors of
short-horizon volatility (absolute-return magnitude) at 5\,s. The
family is entirely different in character: self-excitation of trade
arrivals (a Hawkes-type intensity), realized and range-based
volatility over trailing windows, raw activity counts, and an
effective-spread toxicity proxy in the spirit of the flow-toxicity
literature \citep{easley2012}.

\begin{table}[htbp]
\centering
\small
\caption{Leading volatility features: rank IC vs.\ future absolute
return at 5\,s (census window).}
\label{tab:volatility}
\begin{tabular}{lrrr}
\toprule
Feature & BTC & ETH & SOL \\
\midrule
\feat{hawkes\_intensity}       & $+0.345$ & $+0.281$ & $+0.247$ \\
\feat{vol\_returns\_5m}        & $+0.328$ & $+0.326$ & $+0.293$ \\
\feat{vol\_parkinson\_5m}      & $+0.318$ & $+0.312$ & $+0.281$ \\
\feat{flow\_count\_30s}        & $+0.318$ & $+0.281$ & $+0.253$ \\
\feat{illiq\_trade\_count}     & $+0.318$ & $+0.286$ & $+0.261$ \\
\feat{toxic\_effective\_spread}& $+0.291$ & $+0.279$ & $+0.254$ \\
\bottomrule
\end{tabular}
\end{table}

Notably, VPIN-style measures and tick-sequence entropy carry
magnitude information only --- they have no directional IC at any
horizon tested, on any symbol.

% ==================================================================
\section{Orthogonality: The Factorization Result}
\label{sec:orthogonality}
% ==================================================================

The paper's central positive finding is the clean separation between
the two channels. Every feature in Table~\ref{tab:directional} has
volatility-IC indistinguishable from zero, and every feature in
Table~\ref{tab:volatility} has directional IC indistinguishable from
zero --- consistently across BTC, ETH and SOL. Direction and
volatility of short-horizon returns are predictable \emph{from
disjoint measurements of the same order book}.

Three consequences follow. First, architecturally: a forecasting
system can be built as two independent blocks --- a sign block driven
by the imbalance family and a size/veto block driven by the intensity
family --- with no interaction terms to estimate and hence no
interaction terms to overfit. Second, econometrically: the
factorization licenses separate estimation and separate validation of
the two blocks, halving the effective multiple-testing burden. Third,
economically: it suggests the two channels reflect different
underlying processes --- resting-liquidity positioning (direction)
versus aggressive-flow clustering (magnitude) --- rather than a single
latent state observed twice. \pending{formal independence tests
(distance correlation / HSIC) beyond IC-level orthogonality, on the
replication month.}

% ==================================================================
\section{Spectral Localization and Efficiency Drift}
\label{sec:spectral}
% ==================================================================

\paragraph{Localization.} Frequency-domain analysis of the level-1
imbalance process shows a red spectrum with slope $-1.86$
(Hurst exponent $H = 0.43$, mildly anti-persistent). The predictive
content is not spread across the spectrum: essentially the
\emph{entire} directional IC is carried by the 0.005--0.1\,Hz band
(periods of 10--200\,s). Fitting an Ornstein--Uhlenbeck process to the
imbalance yields a half-life of 5--7\,s, and the dominant
cross-coherence between imbalance and subsequent returns sits at
${\sim}68$\,s. A practical corollary that we verified directly:
time-domain smoothing (exponentially weighted averaging) of the
imbalance \emph{reduces} IC --- naive aggregation destroys exactly the
band that predicts.

\paragraph{Drift.} Between the May census window (05-19/21) and a
June window (06-02/04), the 5\,s IC on BTC declined from $+0.45$ to
$+0.29$ --- on all symbols in the same direction --- while the 1\,s IC
was stable or improved (BTC $+0.447 \to +0.502$). Over the same
period, quoted spreads tightened from ${\sim}0.1$ to 0.2--0.3\,bps
levels shifted regime. The signal is not disappearing; its
\emph{half-life is shortening}: price discovery is accelerating as the
venue matures. To our knowledge, real-time measurement of a
microstructure signal's decay toward efficiency in a young venue has
not been reported at this granularity. \pending{third time-point from
the July replication month.}

% ==================================================================
\section{Robustness}\label{sec:robustness}
% ==================================================================

Per-day rolling IC at 5\,s is $+0.433 \pm 0.061$ (BTC),
$+0.400 \pm 0.098$ (ETH), $+0.355 \pm 0.071$ (SOL); the bootstrap 95\%
confidence interval for \feat{imbalance\_qty\_l1} on BTC is
$[0.442, 0.465]$ (1{,}000 resamples). The signal is present around the
clock with no dead zones (weakest 12--16 UTC, strongest 04--08 UTC)
and in both volatility regimes under a median split (low-vol
$+0.458/+0.435/+0.391$ vs.\ high-vol $+0.422/+0.394/+0.342$ across
BTC/ETH/SOL). Feature rank ordering is preserved across symbols and
subwindows.

We report the validation matrix honestly: BTC passes 7 of 8 cells,
ETH and SOL 6 of 8. The systematic failure is temporal out-of-sample
drift (IC delta $-0.17$ across a six-week gap, against a $0.10$
tolerance) --- partially attributable to the efficiency drift of
Section~\ref{sec:spectral}, and the direct motivation for the
pre-registered replication in Section~\ref{sec:replication}.

% ==================================================================
\section{The Execution Boundary}\label{sec:boundary}
% ==================================================================

The census ICs are mid-price statements. The economically relevant
question is what survives at the point of execution. Two boundaries
are quantified.

\paragraph{Taker arithmetic.} The expected 1--5\,s move conditional on
a strong signal is 0.5--2\,bps; round-trip taker cost on the venue is
${\sim}11$\,bps. Aggressive capture of the short-horizon channel is
arithmetically impossible, independent of modeling choices.

\paragraph{Maker fills and adverse selection.} The subtler boundary
concerns passive execution. Table~\ref{tab:conditional} conditions
the 5\,s IC of \feat{imbalance\_qty\_l1} on fill events of a resting
order at the touch (mid-cross fill model, 23 gap-audited dates,
2026-04-19 to 06-04).

\begin{table}[htbp]
\centering
\small
\caption{Fill-conditioning decomposition: rank IC of
\feat{imbalance\_qty\_l1} at 5\,s under fill conditioning.}
\label{tab:conditional}
\begin{tabular}{lrrr}
\toprule
Conditioning & BTC & ETH & SOL \\
\midrule
Unconditional            & $+0.453$ & $+0.438$ & $+0.412$ \\
Any fill event           & $+0.526$ & $+0.516$ & $+0.460$ \\
Buy fill (signal long)   & $+0.032$ & $-0.061$ & $-0.032$ \\
Sell fill (signal short) & $-0.047$ & $-0.027$ & $-0.021$ \\
\bottomrule
\end{tabular}
\end{table}

The decomposition is the point. Conditioning on the \emph{presence} of
a fill raises the IC above its unconditional level (ticks near fills
carry larger imbalance). Conditioning on the \emph{directionally
consistent} fill --- the one a signal-following market maker actually
receives --- annihilates it, on all three symbols. The mechanism: the
signal says buy; the bid is placed; the fill requires the mid to fall
to the bid, an adverse move; at fill time the prediction is already
spent. This is adverse selection in the classical sense
\citep{glosten1985,kyle1985}, made visible as a conditional-IC
statement, and it locates the monetization loss at \emph{fill
selection}, not at signal decay. Fill prevalence rises from 20.8\%
(BTC) to 30.8\% (SOL) of ticks, and profitability worsens in the same
order: more fills mean more adverse selection, not more edge. The
result complements \citet{collindufresne2015}: information plainly
visible in the book need not be extractable through the act of trading
on it.

\paragraph{Partial recovery by regime conditioning.} Conditioning on a
book-shape entropy measure recovers part of the boundary: in the
low-entropy (structured-book) quintile, the imbalance IC rises from
$0.45$ to $0.55$ (single condition) and to $0.63$--$0.67$ in
combination --- a $+22\%$ lift replicated across symbols. Structured
books appear to correspond to states where resting-order information
is less contested. \pending{fill-conditional version of the entropy
gate on the replication month --- the decisive test of whether regime
conditioning survives at execution.}

% ==================================================================
\section{Out-of-Time Replication (Pre-Registered)}
\label{sec:replication}
% ==================================================================

At the time of writing, an independent month of data (July 2026) is in
final collection. The following analyses are pre-registered and will
replace this section in the next version:

\begin{enumerate}
\item Full census re-estimation (Tables~\ref{tab:directional},
\ref{tab:volatility}, \ref{tab:decay}) on the new month; the
temporal-OOS validation cell re-scored.
\item Fill-conditioning decomposition (Table~\ref{tab:conditional})
re-estimated on the independent window.
\item Orthogonality re-tested, including formal independence measures.
\item A third time-point for the efficiency-drift trajectory.
\item Entropy-gate lift re-estimated unconditionally and
fill-conditionally.
\end{enumerate}

All tests run under Benjamini--Hochberg false-discovery control across
the full battery, with embargoed walk-forward splits where model
fitting is involved; transaction-cost parameters come from a single
audited configuration source. Kill criteria are stated in advance: a
claim that fails replication is removed, not re-tuned.

% ==================================================================
\section{Implications and Future Work}\label{sec:implications}
% ==================================================================

\paragraph{Two-block architectures.} The factorization of
Section~\ref{sec:orthogonality} suggests forecasting and trading
systems in which a directional block (imbalance family) sets the sign
and a volatility block (intensity family) sets the size and the veto
--- e.g.\ reservation-price skew and spread width in an
Avellaneda--Stoikov market maker \citep{avellaneda2008,gueant2013},
with quoting disarmed when toxicity predictors fire
\citep{cartea2018}. The microprice of \citet{stoikov2018} is the
natural implementation of the sign block at the quote level.

\paragraph{The queue-priority question.} The fill-conditioning
collapse of Section~\ref{sec:boundary} is a property of
\emph{reactive} limit orders. Whether orders placed ex ante ---
holding queue priority, whose fills sample noise flow rather than only
adverse flow \citep{moallemi2016} --- retain positive fill-time IC is,
in our view, the central open question this measurement program poses.
A queue-position fill simulation on recorded raw trade data is in
progress.

\paragraph{Execution overlays.} Even absent standalone monetization,
a 0.45-IC seconds-horizon predictor can time the entries of slower,
independently-validated strategies, converting forecast into
execution-cost reduction that fill-selection effects cannot reach
(the trader chooses when to cross, never resting). Quantifying the
recovered basis points per trade is left to companion work.

\paragraph{Efficiency dynamics.} The drift result invites a
market-design reading: a young venue's order book becoming
informationally efficient in measurable real time. Longer histories
will support formal tests of convergence rates.

% ==================================================================
\section{Conclusion}\label{sec:conclusion}
% ==================================================================

The perpetual-futures order book carries two large, orthogonal
predictive channels: a directional channel led by book imbalance
(IC up to $0.47$ at seconds horizons, universal across symbols,
spectrally localized at 10--200\,s periods) and a volatility channel
led by arrival intensity (IC up to $0.35$), with no measurable
cross-talk between the two families. The directional channel is
bounded at execution by adverse selection --- visible as a
fill-conditioning decomposition that raises IC on fill presence and
annihilates it on fill direction --- and partially recovered by
book-entropy regime conditioning. The signal's half-life is shortening
as the venue matures. What the book knows, where it knows it, and
where that knowledge stops being capturable can all be measured; this
paper is an attempt to do so honestly.

% ==================================================================
% REFERENCES
% ==================================================================
\begin{thebibliography}{99}
\setlength{\itemsep}{2pt}

\bibitem[Ackerer et al.(2025)]{ackerer2025}
Ackerer, D., Hugonnier, J., and Jermann, U. (2025).
Perpetual futures pricing. \emph{Mathematical Finance}, forthcoming.

\bibitem[Avellaneda and Stoikov(2008)]{avellaneda2008}
Avellaneda, M., and Stoikov, S. (2008).
High-frequency trading in a limit order book.
\emph{Quantitative Finance}, 8(3), 217--224.

\bibitem[Benjamini and Hochberg(1995)]{benjamini1995}
Benjamini, Y., and Hochberg, Y. (1995).
Controlling the false discovery rate: a practical and powerful
approach to multiple testing.
\emph{Journal of the Royal Statistical Society B}, 57(1), 289--300.

\bibitem[Cartea et al.(2018)]{cartea2018}
Cartea, \'A., Donnelly, R., and Jaimungal, S. (2018).
Enhancing trading strategies with order book signals.
\emph{Applied Mathematical Finance}, 25(1), 1--35.

\bibitem[Collin-Dufresne and Fos(2015)]{collindufresne2015}
Collin-Dufresne, P., and Fos, V. (2015).
Do prices reveal the presence of informed trading?
\emph{Journal of Finance}, 70(4), 1555--1582.

\bibitem[Cont et al.(2014)]{cont2014}
Cont, R., Kukanov, A., and Stoikov, S. (2014).
The price impact of order book events.
\emph{Journal of Financial Econometrics}, 12(1), 47--88.

\bibitem[Easley et al.(2012)]{easley2012}
Easley, D., L\'opez de Prado, M., and O'Hara, M. (2012).
Flow toxicity and liquidity in a high-frequency world.
\emph{Review of Financial Studies}, 25(5), 1457--1493.

\bibitem[Glosten and Milgrom(1985)]{glosten1985}
Glosten, L., and Milgrom, P. (1985).
Bid, ask and transaction prices in a specialist market with
heterogeneously informed traders.
\emph{Journal of Financial Economics}, 14(1), 71--100.

\bibitem[Gu\'eant et al.(2013)]{gueant2013}
Gu\'eant, O., Lehalle, C.-A., and Fernandez-Tapia, J. (2013).
Dealing with the inventory risk: a solution to the market making
problem. \emph{Mathematics and Financial Economics}, 7(4), 477--507.

\bibitem[Kyle(1985)]{kyle1985}
Kyle, A. (1985).
Continuous auctions and insider trading.
\emph{Econometrica}, 53(6), 1315--1335.

\bibitem[Moallemi and Yuan(2016)]{moallemi2016}
Moallemi, C., and Yuan, K. (2016).
The value of queue position in a limit order book.
Working paper, Columbia Business School.

\bibitem[Stoikov(2018)]{stoikov2018}
Stoikov, S. (2018).
The micro-price: a high-frequency estimator of future prices.
\emph{Quantitative Finance}, 18(12), 1959--1966.

\end{thebibliography}

% ==================================================================
\appendix
\section{Reproducibility}\label{app:repro}
% ==================================================================

\begin{itemize}
\item \textbf{Code revision:} git SHA \texttt{[stamp at submission]}
of the research platform (Rust ingestion engine + Python analysis).
\item \textbf{Census window:} 2026-05-19 to 2026-05-21;
2.17\,M ticks/symbol at 100\,ms cadence; 21{,}708 subsampled IC
points per symbol; horizons 1\,s--15\,m.
\item \textbf{Fill-conditioning window:} 2026-04-19 to 2026-06-04;
23 of 31 dates pass gap audit; mid-cross fill model at the touch.
\item \textbf{Drift windows:} 2026-05-19/21 vs.\ 2026-06-02/04.
\item \textbf{Replication month:} July 2026 (in collection at draft
time); pre-registered battery in Section~\ref{sec:replication}.
\item \textbf{Multiplicity:} Benjamini--Hochberg FDR across each
reported battery; significance $p<0.01$ for census entries.
\item \textbf{Costs:} all cost statements from a single audited fee
configuration (taker round-trip ${\sim}11$\,bps at study time).
\end{itemize}

Data samples and analysis scripts are available from the author on
request.

\end{document}
